\documentclass[a4paper,11pt]{ctexart}
\usepackage{xcolor}
\usepackage{array}
\usepackage{booktabs}
\usepackage{hyperref}
\hypersetup{colorlinks=true,linkcolor=blue!70!black}
\linespread{1.35}

\newcommand{\draft}{\textcolor{orange!80!black}{\textbf{【待写】}}}

\title{\textcolor{blue!70!black}{\Huge\textbf{航空航天模块 · 课程大纲}}}
\author{}
\date{}

\begin{document}
\maketitle
\thispagestyle{empty}

\section*{第零层 · 入门故事}

\begin{enumerate}
\item[\textbf{第1课}] \textbf{火箭是怎么飞上天的？——牛顿第三定律的极致应用} \draft \\
\textbf{内容：} 向下喷气→向上飞（作用力与反作用力）。和气球放气的原理完全一样。多级火箭的奥秘（一节烧完丢掉——越来越轻——越来越快）。长征五号和天宫空间站的故事。
\item[\textbf{第2课}] \textbf{宇航员在太空怎么生活？——失重不是没有引力} \draft \\
\textbf{内容：} 空间站在绕地球"自由落体"——宇航员和站体一起"掉"，感觉不到重量。在太空吃饭、睡觉、上厕所的趣事。为什么宇航员必须每天锻炼两小时？
\item[\textbf{第3课}] \textbf{飞机为什么能飞起来？——伯努利原理与反作用力} \draft \\
\textbf{内容：} 机翼上表面空气流速快、压强低；下表面流速慢、压强高——产生升力。但不要忽略"攻击角"也很重要。纸飞机飞行实验——什么形状飞得最远？
\end{enumerate}

\section*{深入课程}

\begin{enumerate}
\item[\textbf{第1课}] \textbf{空气动力学——飞机如何对抗重力和阻力} \draft \\
\textbf{内容：} 升力公式 $L = \frac12\rho v^2 S C_L$、阻力公式、升阻比。翼型设计、高升力装置（襟翼/前缘缝翼）。激波与音障。情景：为什么客机巡航在10000米高度而不是3000米？
\item[\textbf{第2课}] \textbf{推进系统——从活塞发动机到离子推进器} \draft \\
\textbf{内容：} 涡轮风扇发动机工作原理（进气道→压气机→燃烧室→涡轮→喷管）、推力公式 $F = \dot{m}(v_e - v_0)$。比冲、推重比。火箭发动机的泵压式 vs 挤压式。情景：为什么火箭要用多级？（Tsiolkovsky火箭方程$\Delta v = v_e \ln(m_0/m_f)$）
\item[\textbf{第3课}] \textbf{轨道力学——卫星为什么不掉下来？} \draft \\
\textbf{内容：} 牛顿的"高山大炮"思想实验→轨道速度。开普勒三定律、轨道根数。GTO/LEO/GEO/SSO 轨道的区别和应用。霍尔曼转移轨道。情景：为什么SpaceX的星链卫星要在550km高度？（不是更高更低，而是考虑了大气阻力和范艾伦辐射带）
\item[\textbf{第4课}] \textbf{航天器设计——在太空里活下去} \draft \\
\textbf{内容：} 热控（太空中向阳面150°C、背阴面-150°C）、姿控（飞轮/推力器/磁力矩器）、电源（太阳能+蓄电池）、通信（深空通信的链路预算）。情景：旅行者1号发射已47年、距离地球240亿公里——它是怎么还能发回信号的？
\item[\textbf{第5课}] \textbf{制导与导航——如何让火箭精确入轨} \draft \\
\textbf{内容：} 惯性导航（IMU）+ GPS/北斗修正。火箭飞行的"程序转弯"。再入大气层的"黑障"（等离子体切断通信）。情景：嫦娥五号从月球采样返回——再入大气层时的速度是11km/s——如何确保不烧毁？
\item[\textbf{第6课}] \textbf{北京在地——中国航天和北京} \draft \\
\textbf{内容：} 中国运载火箭技术研究院（南苑）、中科院国家空间科学中心（怀柔）、北京航空航天大学（学院路）。参观指南：中国航天博物馆（丰台）、北京天文馆、北航航空航天博物馆（学院路校区）。
\end{enumerate}

\end{document}
